% ==========================================================================
% GENERATED by scripts/split_paper.py. Do not edit this file.
%
% Source: imported_papers/Neural_Computation__CNA/NECO-D-24-00304-Source.tex
% Part 2: The Net Fragment Framework
%
% Edits belong upstream, in the published paper. To change what the
% chapter includes, change imports/cna.toml and re-run
% the splitter.
% ==========================================================================
\section{The Net Fragment Framework}
\label{sec:fundamentals}
Before presenting the computational framework, we introduce its foundational components: neurons, activations, nets, net fragments, and competition. These theoretical definitions serve as the basis for the subsequent sections. Implementation details for the introduced formalisms are provided in the next section.


\paragraph{Neurons.}
A neuron $i$ is denoted by $n_i$. Each neuron $n_i$ has a \emph{local neighborhood}, denoted by the closed ball $\mathcal{B}_r(n_i)$, which is the set of neurons $n_j$ such that the distance $d(n_i, n_j)$ between $n_i$ and $n_j$ is less than or equal to a threshold radius $r>0$, i.e.,
%
\begin{equation}
\mathcal{B}_r\left(n_i\right) = \left \{n_j \mid d(n_i, n_j) \leq r \right \}.
\label{eq:neighborhood}
\end{equation}
%
At each time step $t$, the \emph{activity state} of a neuron $n_i$ is represented by $y_{i,t} \in \{0, 1\}$, where $y_{i,t} = 1$ indicates that the neuron is active at $t$. The activity of a neuron over time is characterized by an \emph{activation pattern}, which is a sequence of activation states $\{y_{i,t}\}_{t=1}^T$, where $T$ is the total number of time steps (sequential processing steps).

\paragraph{Connections.}
Neurons interact via \emph{connections}. A direct connection from neuron $n_i$ to another neuron $n_j$ is defined by a positive weight $w_{i,j} > 0$, which facilitates the propagation of neuronal activity from $n_i$ to $n_j$, so that the activity $y_{i,t}$ of neuron $n_i$ can influence the activity $y_{j,t+1}$ of neuron $n_j$ at the subsequent time step. In this work, we limit the connections to local neighborhoods:
%
\begin{equation}
\forall  w_{i,j} \quad n_j \in \mathcal{B}_r(n_i)
\label{eq:neighborhood_connections}
\end{equation}
%
Additionally, we limit the weights to a range between $0$ and $1$, allowing neuron $n_i$ only to positively influence $n_j$:%
\begin{equation}
w_{i,j} \in [0,1]
\label{eq:neighborhood_connections_range}
\end{equation}
%
We define a relation $n_i \sim n_j$ between two neurons to indicate that their activation patterns are highly correlated, i.e., $y_{i,t}$ often influences  $y_{j,t+1}$.
Under Hebbian learning \citep{hebb_organization_1949}, the connection strength $w_{i,j}$ increases when there is $n_i \sim n_j$.


\paragraph{Net Fragments.}
A \emph{net fragment}  is defined as the set of neurons that provide direct activation support to a given neuron.
A neuron $n_j$ is said to provide \emph{support} to another neuron $n_i$ if two conditions are met:
(1) there exists a high correlation between two neurons $n_j$ and $n_i$ such that $n_j \sim n_i \Rightarrow w_{ji} > 0$; and (2) neuron $n_j$ was active at the previous time step, i.e., $y_{j,t-1}=1$.
Since synaptic weights are bounded within the interval $w_{j,i} \in [0,1]$, the degree of support a neuron can provide is likewise restricted to this range.

Formally, the \emph{net fragment} associated with neuron $n_i$ at time $t$ is given by
%
\begin{equation}
  \mathcal{F}_{t}(n_i) = \{n_i\} \cup \{n_j \mid w_{j,i} > 0 \land y_{j,t-1} = 1\}.
  \label{eq:net_fragments}
\end{equation}
%
In this formulation, any neuron $n_j \in \mathcal{F}(n_i)$  contributes to the activation likelihood of $n_i$ at time $t$.
Notably, neuron $n_i$ is part of its own fragment, reflecting a form of self-support -- if active at time $t$, the neuron may contribute to sustaining its own activation at the subsequent time step.
Note that $n_i \sim n_j$ is not an equivalence relation since correlation is not transitive, meaning that one neuron can be part of multiple net fragments.

\paragraph{Activation Function.}
The decision of whether a neuron becomes active is based on the amount of support it receives from within the net fragment. We call this accumulated input the \emph{pre-synaptic activity} at timestep $t$, which is a function $\mathcal{I}$ of the net fragment $\mathcal{F}_{t-1}(n_i)$ of a neuron $n_i$ at timestep $t-1$. %that is calculated based on the neurons' state that are part of its fragment:
%
\begin{equation}
  \boldsymbol{a}^\mathcal{I}_{i, t} = \mathcal{I}\left(\mathcal{F}_{t}(n_i) \right)  
  \label{eq:pre_synaptic_activity}
\end{equation}
%
The pre-synaptic activity $a^\mathcal{I}_{i, t} \in \mathbb{R}$ is large when many other neurons provide support.
Since supporting neurons exhibit connections that have been selectively strengthened through previous mutual activity and Hebbian updates within specific contexts, their activity effectively serves to confirm the contextual plausibility of a given neuron's activation.
The very general formulation of $\mathcal{I}$ shows that the pre-synaptic activity can be calculated with any desired function.
In this work, we specifically define $\mathcal{I}$ as the weighted sum of prior activations from supporting neurons:
%
\begin{equation}
\boldsymbol{a}^\mathcal{I}_{i, t} = \sum_{n_j \in \mathcal{F}_{t}(n_i)} \boldsymbol{w}_{j,i} \cdot \boldsymbol{y}_{j,t-1}.
\label{eq:specific_pre_synaptic_activity}
\end{equation}


The activation of a neuron at timestep $t$ is determined by the activation function $\mathcal{A}$ that depends on the pre-synaptic activity $a_{i, t}$ and environmental conditions $\mathcal{E}$: 
%To ensure that only well-supported neurons remain active, they undergo inhibition, whose strength is determined by the environment denoted as $env$, and are then binarized by an activation function:
%
\begin{equation}
  \boldsymbol{y}_{i, t} = \mathcal{A} \left(\boldsymbol{a}_{i, t}, \mathcal{E} \right) \in \{0,1\}
  \label{eq:activation_framework}
\end{equation}
%
The environmental conditions $\mathcal{E}$ include an \emph{attenuation} signal, a constraint that selectively suppresses weak excitatory inputs by requiring neurons to accumulate a sufficiently large pre-synaptic activity before firing.
Again, we here provide a very general formulation, using $\mathcal{A}$ and $\mathcal{E}$, to emphasize that any activation function or environmental mechanism can be used.
In this work, we instantiate $\mathcal{E}$ using a combination of normalization (see Appendix~\ref{sec:normalization}) and attenuation (see Eq.~\ref{eq:inhibition}), and define $\mathcal{A}$ as the heaviside step function (see Eq.~\ref{eq:s2_final_output}). 

\paragraph{Nets.}
Neurons are not only influenced by direct connections within net fragments but also by indirect connections (e.g., if $n_i$ supports $n_j$ and $n_j$ supports $n_k$, then $n_i$ indirectly supports $n_k$).
A set of directly and indirectly supported neurons is called a \emph{net}, implementing the basic structure used in this work to represent objects.
For simplicity in the following definitions, we introduce $\mathcal{F}^*_{t}(n_i)$, a net fragment $\mathcal{F}_{t}(n_i)$ in which $n_i$ is active as well:
%
\begin{equation}
  \mathcal{F}^*_{t}(n_i) =\mathcal{F}_{t}(n_i) \quad \text{with} \quad  y_{i, t}=1
  \label{eq:net_fragments_full}
\end{equation}
%
To account for both direct and indirect influences on a neuron's activation, we use a recursive formulation of nets from the perspective of a single neuron. In the base case, a net $\mathcal{N}_t^{1}[n_i]$ is defined as the immediate net fragment:
%
\begin{equation}
    \mathcal{N}_t^{1}[n_i]  = \mathcal{F}^*_t(n_i)
\end{equation}
%
For $r \geq 2$, the recursive relation is given by
%
\begin{equation}
    \mathcal{N}_t^{r}[n_i]  = \mathcal{N}_t^{r-1}[n_i] \bigcup_{n_j \in \mathcal{N}_t^{r-1}[n_i]} \mathcal{F}^*_t(n_j)
    \label{eq:recursive_net_fragments}
\end{equation}
%
The complete network accessible from $n_i$ is then defined as the closure
%
\begin{equation}
    \mathcal{N}_t^{*}[n_i]  = \bigcup_{r=1}^{\infty} \mathcal{N}_t^{r}[n_i],
\end{equation}
%
which comprises all neurons in the active subgraph that are reachable from $n_i$.
This recursive formulation illustrates how support relationships propagate through time: local net fragments $\mathcal{F}_t(n_i)$, which reflect immediate sources of support for a neuron, are combined recursively to form increasingly larger subnetworks $\mathcal{N}_t^{r}[n_i]$. These progressively include neurons with indirect but recurrent influence on $n_i$'s activation. This process reflects how localized evidence (elementary activations) can aggregate into a coherent larger pattern (composite activation), corresponding to a higher-order feature representation in the network.
Moreover, the explicit time dependence in these definitions underscores the dynamic nature of both the net fragments and the overall net, as they are continuously reassembled in response to the evolving support received from neighboring neurons, and increasing attenuation that prevents neurons with a low pre-synaptic activity from firing (see next section).

In the proposed architecture, no symmetry constraint is imposed on the weights, that is, $w_{i,j} \neq w_{j,i}$ in general.
However, due to the temporal dynamics of the network, activations tend to stabilize over a longer time horizon $T$, such that for sufficiently large $t$, the activations become approximately stationary, i.e., $y_{i, t} \approx y_{i, t+1}$.
Under these conditions, the mutual influence between neurons becomes symmetric, meaning that $n_i \sim n_j$ implies $n_j \sim n_i$, and almost symmetric weights emerge naturally over time, even in the absence of explicit symmetry constraints.


As a result, it can be assumed that for any pair of distinct neurons $n_i \neq n_j$, their closures either converge to the same set of neurons or remain largely disjoint.
Formally, this can be expressed as:
%
\begin{equation}
    \mathcal{N}_t^{*}[n_i]  \approx \mathcal{N}_t^{*}[n_j] \quad \text{or} \quad  \mathcal{N}_t^{*}[n_i]  \cap  \mathcal{N}_t^{*}[n_j] \approx \emptyset.
\end{equation}


\paragraph{Competitive Neurons.}
One issue with the current formulation is that the mapping
%
\begin{equation}
  \mathcal{A}\circ\mathcal{I}:2^{\mathcal{N}}\;\longrightarrow\;\{0,1\}
\end{equation}
%
treats any two net fragments that yield the same activation identically.  In particular, we say
%
\begin{equation}
  \mathcal{F}_{t_k}(n_i)\simeq\mathcal{F}_{t_l}(n_i)
  \quad\Longleftrightarrow\quad
  \mathcal{A}\bigl(\mathcal{I}(\mathcal{F}_{t_k}(n_i)),\mathcal{E}\bigr)
  =\mathcal{A}\bigl(\mathcal{I}(\mathcal{F}_{t_l}(n_i)),\mathcal{E}\bigr),
\end{equation}
%
meaning different contexts collapse to the same output for $n_i$. To break this symmetry and make each neuron context-aware, we introduce \emph{competitive neurons}.

For each neuron $n_i$, We introduce a set of alternatives $\{n_i^1,\dots,n_i^\kappa\}$, all located at the same spatial position $\left ( \mathcal{B}_r(n_i)=\mathcal{B}_r(n_i^k) \right )$.
These alternatives undergo competition and one of them is permitted to be active within a given context. This exclusivity encourages these competitive neurons to specialize, over time, into representations of different contextual variations.

To formalize these competitive neurons, we first define the \emph{competition class} of a neuron $n_i$:
%
\begin{equation}
  \mathcal{C}(n_i) 
  = \bigl\{ n_i^k  \, \text{ where } \, \mathcal{B}_r(n_i)=\mathcal{B}_r(n_i^k) \bigr\},
\end{equation}
%
These neurons are mutually exclusive candidates to represent the context-dependent expression of a feature at a fixed spatial position.
During each time step, after computing the pre-synaptic activities $\boldsymbol{a}_{j,t}^\mathcal{I}$ for $n_j\in\mathcal{C}(n_i)$, we select a small winner set
%
\begin{equation}
  W_t(n_i)
  = \mathrm{Comp}\bigl(\{\,\boldsymbol{a}_{j,t}^\mathcal{I}\mid n_j\in\mathcal{C}(n_i)\}\bigr).
\end{equation}
%
While various rules to determine the winning set can be used, we apply winner-takes-all selection in this work (for implementation details, see Section~\ref{sec:stage_2_impl}).
Only neurons in $W_t(n_i)$ are eligible to fire:
%
\begin{equation}
  y_{i,t} \;=\;
  \begin{cases}
    \mathcal{A}\bigl(\boldsymbol{a}_{i,t}^\mathcal{I},\mathcal{E}\bigr)
      & \text{if }n_i\in W_t(n_i),\\
    0 & \text{otherwise.}
  \end{cases}
\end{equation}

By enforcing mutual exclusivity among neurons with overlapping receptive fields, competition ensures that different contexts activate distinct neuronal subsets.
